% Paper 3, §2 — The CoT phenomenon and its scale-emergence
% Anchors: kb/notes/reasoning/chain-of-thought.md §1-2
%          kb/excerpts/{wei2022-cot,kojima2022}.md
% Forward refs: §3 self-consistency (sec:self-consistency),
%               §4 inference-time compute (sec:test-time-compute),
%               §11 faithfulness (sec:cot-faithfulness)

\section{The CoT phenomenon and its scale-emergence}
\label{sec:cot-phenomenon}

The historical entry point of every leg of the reasoning triangle is a
single decoding pattern: instead of mapping a problem $x$ directly to an
answer $y$, the model first produces an intermediate \emph{reasoning
trace} $z$ and only then commits to $y$. \citet{wei2022-cot} introduced
this pattern as a few-shot prompting recipe under the label
\textbf{\glsterm{chain-of-thought}{chain-of-thought} (\glsterm{chain-of-thought}{CoT})} prompting, with the headline finding
that an 8-exemplar \glsterm{chain-of-thought}{CoT} prompt raises a 540B-parameter language model on
the GSM8K math word-problem benchmark to state-of-the-art, surpassing a
fine-tuned GPT-3 with a verifier
\citep[abstract]{wei2022-cot}.\footnote{KB excerpt:
\texttt{kb/excerpts/wei2022-\glsterm{chain-of-thought}{cot}.md\#abstract}.} The same pattern was
shown a few months later to elicit competitive reasoning under
\emph{zero} hand-written exemplars, by appending the literal trigger
phrase \emph{``Let's think step by step''} before answer extraction
\citep[abstract]{kojima2022}.\footnote{KB excerpt:
\texttt{kb/excerpts/kojima2022.md\#abstract}.} The 2022 community read
both results as prompting tricks. The 2024--2026 reframing — anchored
in \citet{snell2024} and the inference-time-compute scaling literature
that \cref{sec:test-time-compute} formalises — reads them as the first
empirical instance of \emph{trading inference compute for \glsterm{parameters}{parameters}}.
This section walks the original mechanism, transcribes the decoding
factorisation that every later test-time-compute strategy in the paper
will build on, presents the emergence-with-scale finding that
distinguished \glsterm{chain-of-thought}{CoT} from compute-monotonic prompting, and closes by
deferring \glsterm{chain-of-thought}{CoT}'s faithfulness question to \cref{sec:cot-faithfulness}.

\subsection{The decoding factorisation}
\label{sec:cot-decoding}

Let $x$ denote the problem and $y$ the desired final answer. Direct
prompting decodes $y \sim p_\theta(\cdot \mid x)$ in a single step. \glsterm{chain-of-thought}{CoT}
prompting introduces an intermediate variable $z$ — the reasoning trace
— and decodes a two-stage factorisation
\citep[\S 3.1]{wei2022-cot}:\footnote{KB excerpt:
\texttt{kb/excerpts/wei2022-\glsterm{chain-of-thought}{cot}.md\#sec-3-1}.}
\begin{equation}
  z \;\sim\; p_\theta(\cdot \mid x), \qquad y \;\sim\; p_\theta(\cdot \mid x, z).
  \label{eq:cot-decode}
\end{equation}
In practice the trace and answer are emitted jointly by \glsterm{autoregressive}{autoregressive}
decoding of a single concatenated sequence $z \cdot y$, where $z$ is a
natural-language explanation terminating in a stable answer-extraction
template (e.g.\ \emph{``The answer is~$\langle\cdot\rangle$''}) and $y$
is recovered post-hoc by a regex on that template
\citep[\S 3.1]{wei2022-cot}. The KB note for this paper records the
factorisation as Eq.~1 of \texttt{\glsterm{chain-of-thought}{chain-of-thought}.md}.\footnote{KB
note: \texttt{kb/notes/reasoning/\glsterm{chain-of-thought}{chain-of-thought}.md\#1}.}

\paragraph{Tensor shapes and the trace-length budget.} The trace $z$ in
\cref{eq:cot-decode} is a token sequence whose length empirically spans
two orders of magnitude across protocols and benchmarks: the original
GSM8K few-shot \glsterm{chain-of-thought}{CoT} exemplars are on the order of $10^2$ tokens per
trace \citep[\S 3.1]{wei2022-cot}, while the long-\glsterm{chain-of-thought}{CoT} reasoning models
of \cref{sec:r1-family} routinely produce $10^3$--$10^4$ tokens of $z$
per problem \citep[\S 2.3]{deepseek-r1}. The answer $y$ remains the
same shape it was under direct prompting (a few tokens). The
implication for accounting is structural: under \cref{eq:cot-decode}
the inference-time compute and KV-cache budget are dominated by $z$,
not by $y$. Decoding a single \glsterm{chain-of-thought}{CoT} response of shape
$\mathsf{S} = |z| + |y| \approx |z|$ at hidden dim $\mathsf{D}$ over
$\mathsf{L}$ layers carries a per-layer KV-cache cost
$O(\mathsf{S} \cdot \mathsf{D})$ that is the binding constraint long
before the model's parameter count is. The architectural prerequisites
for tractable long traces — long-context attention, KV compression,
flash attention — are the subject of Paper~1 and re-enter this paper
in \cref{sec:r1-family} as the precondition for $10^4$-token $z$.

\paragraph{Few-shot exemplar count $K$.} The \glsterm{chain-of-thought}{CoT} prompt in
\citet{wei2022-cot} is a sequence of $K$ in-context demonstrations
$\{(x_i, z_i, y_i)\}_{i=1}^K$ followed by the \glsterm{query}{query} $x_*$, with each
$z_i$ a hand-written reasoning trace. The arithmetic-reasoning
benchmarks in \citet{wei2022-cot} use $K = 8$
\citep[\S 3.1]{wei2022-cot}.\footnote{KB excerpt:
\texttt{kb/excerpts/wei2022-\glsterm{chain-of-thought}{cot}.md\#sec-3-1}.} The exemplars do double
duty: they (i) shift the model's output distribution toward sequences
in which $z$ precedes $y$ — the format-conditioning role — and
(ii) supply task-specific reasoning style. Kojima et al.'s zero-shot
result below shows the format-conditioning role does not in fact
require $K > 0$, suggesting the exemplars' marginal contribution at
sufficient scale is style, not format.

\subsection{Few-shot CoT (Wei et al.\ 2022)}
\label{sec:cot-fewshot}

The original recipe is the simplest version of \cref{eq:cot-decode}:
construct a prompt of $K = 8$ exemplars with hand-written reasoning
traces, decode greedily (or at low temperature) from $p_\theta(\cdot
\mid \text{prompt}, x_*)$ until the stop pattern, and parse the
trailing \emph{``The answer is~$\langle\cdot\rangle$''} suffix to
recover $y_*$ \citep[\S 3.1]{wei2022-cot}. \citet{wei2022-cot} apply
this protocol to three model families (LaMDA, GPT-3, \glsterm{palm}{PaLM}) at sizes
spanning roughly $0.4\text{B}$ to $540\text{B}$ \glsterm{parameters}{parameters}, evaluated
on three task families: arithmetic word problems (GSM8K, SVAMP, ASDiv,
AQuA, MAWPS), commonsense (CSQA, StrategyQA, Date, Sports), and
symbolic reasoning (Last Letter, Coin Flip)
\citep[\S 3.2]{wei2022-cot}.\footnote{KB excerpt:
\texttt{kb/excerpts/wei2022-\glsterm{chain-of-thought}{cot}.md\#sec-3-2}.} The headline
\emph{single-number} result is GSM8K at the 540B scale: an 8-exemplar
\glsterm{chain-of-thought}{CoT} prompt achieves state-of-the-art accuracy, surpassing a fine-tuned
GPT-3 supplemented with a verifier
\citep[abstract]{wei2022-cot}.\footnote{KB excerpt:
\texttt{kb/excerpts/wei2022-\glsterm{chain-of-thought}{cot}.md\#abstract}.}

\subsection{Zero-shot CoT (Kojima et al.\ 2022)}
\label{sec:cot-zeroshot}

\citet{kojima2022} report that the entire \glsterm{chain-of-thought}{CoT}-eliciting effect can be
recovered with $K = 0$ — no hand-written exemplars at all — by
appending the literal string \emph{``Let's think step by step''} to the
prompt. The protocol is two-stage
\citep[\S 3]{kojima2022}:\footnote{KB excerpt:
\texttt{kb/excerpts/kojima2022.md\#sec-3}.}
\begin{enumerate}
\item \textbf{Reasoning extraction.} Prompt the model with
$x_* \,+\,$\emph{``Let's think step by step''} and decode the trace
$z_*$.
\item \textbf{Answer extraction.} Prompt the model with
$x_* \,+\, z_* \,+\,$\emph{``Therefore, the answer is''} and decode
$y_*$.
\end{enumerate}
The two-stage form is the canonical protocol; the trigger phrase is
\emph{``Let's think step by step''}, identified empirically among
candidate variants by \citet{kojima2022} as the strongest cue. The
headline numbers \citep[headline]{kojima2022}\footnote{KB excerpt:
\texttt{kb/excerpts/kojima2022.md\#headline}.} are striking for a
purely prompting-side change: with InstructGPT (\texttt{text-davinci-002}),
MultiArith rises from $17.7\%$ to $78.7\%$ and GSM8K from $10.4\%$ to
$40.7\%$, with comparable gains on \glsterm{palm}{PaLM}-540B.

The structural consequence: at sufficient scale the format-conditioning
role of Wei's $K = 8$ exemplars is already supplied by pretraining and
instruction-tuning, and a five-word trigger reaches into that latent
capability without a single demonstration
\citep[\S 5]{kojima2022}.\footnote{KB excerpt:
\texttt{kb/excerpts/kojima2022.md\#sec-5}.} The reasoning content is
not absent in the smaller-prompt setting; it is just not elicited.

\begin{intuition}
\emph{``Let's think step by step''} is a soft retrieval cue, not a
reasoning module. Pretraining text — math tutorials, instructional
material, worked examples — contains many instances of that exact
phrase preceding step-by-step explanations, so appending it shifts
the conditional distribution $p_\theta(z, y \mid x_* \,+\, \text{cue})$
toward sequences in which a trace precedes the answer. Returning to
math: the cue does not change the family
\cref{eq:cot-decode}; it only nudges $p_\theta(\cdot \mid x)$ to put
mass on $z$-then-$y$ continuations rather than on direct-answer
continuations. The mechanism is a learned co-occurrence
\citep[\S 5]{kojima2022}, not a discovered reasoning faculty.
\end{intuition}

\subsection{Emergence with scale}
\label{sec:cot-emergence}

The single most consequential empirical claim of
\citet{wei2022-cot} for the rest of this paper is that the gain from
\glsterm{chain-of-thought}{CoT} prompting over direct prompting is itself an
\emph{emergent} property of model scale. The headline figure
\citep[\S 3.3, fig.\ 4]{wei2022-cot}\footnote{KB excerpt:
\texttt{kb/excerpts/wei2022-\glsterm{chain-of-thought}{cot}.md\#sec-3-3}.} shows GSM8K accuracy
plotted against model size for three families (LaMDA, GPT-3, \glsterm{palm}{PaLM}),
each with and without the 8-exemplar \glsterm{chain-of-thought}{CoT} prompt. Two qualitative
features survive every replication:
\begin{itemize}
\item \textbf{No gain — sometimes a loss — below ${\sim}100\text{B}$
\glsterm{parameters}{parameters}.} Below the inflection, \glsterm{chain-of-thought}{CoT} prompting often
\emph{underperforms} direct prompting on GSM8K. The model emits a
trace whose internal arithmetic is brittle, and the trace's errors
propagate into the answer.
\item \textbf{Substantial gain at \glsterm{palm}{PaLM}-540B.} Above the inflection,
\glsterm{chain-of-thought}{CoT} prompting opens a sharp gap over direct prompting — the gap that
produces the abstract's headline of state-of-the-art GSM8K
\citep[abstract]{wei2022-cot}.
\end{itemize}
\citet[\S 6]{wei2022-cot} explicitly frame this in their limitations
discussion: the technique's benefit is not a uniform multiplicative
factor over model scale but a regime-dependent effect that activates
only past a problem-dependent capacity threshold.\footnote{KB excerpt:
\texttt{kb/excerpts/wei2022-\glsterm{chain-of-thought}{cot}.md\#sec-6}.} The same paper hedges
that ``\glsterm{chain-of-thought}{chain-of-thought}'' need not reflect any cognitive sense of
``reasoning'' — a hedge that the 2025 faithfulness literature
(\cref{sec:cot-faithfulness}) makes load-bearing.

\subsection{The post-2024 reframing: prompting trick vs.\ compute axis}
\label{sec:cot-reframing}

Reading the 2022--2023 reception of \glsterm{chain-of-thought}{CoT} in retrospect, the field largely
treated \cref{eq:cot-decode} as a prompting innovation: a clever
demonstration of in-context learning, evaluated relative to
zero-shot or few-shot baselines, on benchmarks chosen to expose
\emph{prompting} differences. The compute axis was implicit but
unaccounted-for. Wei's 8-exemplar \glsterm{chain-of-thought}{CoT} prompt at 540B emits roughly
$10^2$ tokens of $z$ per problem; the direct-prompting baseline emits
roughly $10^0$. The protocols differ by two orders of magnitude in
inference compute per problem, and the GSM8K gap is the joint product
of \emph{(i)} the model's latent capacity to use that compute and
\emph{(ii)} the prompt that elicited the trace.

\citet{snell2024} promote this reading explicitly: \glsterm{test-time-compute}{test-time compute}
becomes a \emph{scaling-law-shaped} axis with its own functional form,
on which \glsterm{chain-of-thought}{CoT} prompting was the first attested operating point. The
emergence-with-scale finding of
\citet[\S 3.3]{wei2022-cot} is then re-readable as the joint shape of
the inference-compute / parameter-count surface — sub-100B models
under \cref{eq:cot-decode} cannot convert the extra inference compute
into accuracy because their per-token reasoning step is too noisy;
\glsterm{palm}{PaLM}-540B can. The 2025--2026 reasoning-trained models of
\cref{sec:r1-family} push the operating point further along the same
axis: $z$ grows from $10^2$ to $10^4$ tokens, the model is RL-trained
to use that budget, and the gain over a $10^0$-token direct answer
becomes the dominant accuracy lever on hard math benchmarks. \glsterm{chain-of-thought}{CoT} was
not magic prompting; it was the first published instance of substituting
inference compute for \glsterm{parameters}{parameters} along an axis the field would later
formalise.

\begin{analogy}
\glsterm{chain-of-thought}{CoT} is to direct prompting what stochastic gradient descent with $K$
inner-loop steps is to a single gradient step: the same model now
spends additional compute per \glsterm{query}{query} before committing to an output,
and the marginal accuracy returned by that compute is the load-bearing
quantity. Returning to math: the analogy collapses to
\cref{eq:cot-decode} — direct prompting collapses $z$ to a length-zero
trace ($p_\theta(y \mid x)$ alone), \glsterm{chain-of-thought}{CoT} prompting decodes a non-trivial
$|z| > 0$, and \cref{sec:test-time-compute} treats $|z|$ together with
the sample count $N$ of \cref{sec:self-consistency} as orthogonal
compute knobs whose joint optimisation is the inference-side
counterpart of training-side scaling laws \citep[\S 1]{snell2024}.
\end{analogy}

\subsection{Forward link to faithfulness}
\label{sec:cot-faith-forward}

A question \cref{eq:cot-decode} raises but does not settle: \emph{does
$z$ in fact cause $y$?} The factorisation is well-defined as a
generative model — the trace is sampled, the answer is sampled
conditional on the trace — but the conditional dependence
$y \sim p_\theta(\cdot \mid x, z)$ is not the same as a causal claim
that perturbations to $z$ propagate to $y$. The 2025--2026 faithfulness
literature \citep{arcuschin2025-cot-wild,shen2025-faithcot-bench,
mehta2026-faithfulness-scaling} measures this directly and finds the
answer is regime-dependent.

\begin{contradiction}
\textbf{Reasoning vs.\ rationalisation.} \cref{eq:cot-decode} treats
$z$ as a sampled latent that conditions $y$, but the operational
question is whether perturbations to $z$ that should change $y$ in fact
do. Production non-thinking models exhibit substantial post-hoc
rationalisation rates (on the order of $7$--$13\%$ on
\citet{mehta2026-faithfulness-scaling}'s measurement); thinking models
exhibit rates two orders of magnitude smaller
\citep[\S 4]{mehta2026-faithfulness-scaling}.\footnote{KB note:
\texttt{kb/notes/reasoning/\glsterm{chain-of-thought}{chain-of-thought}.md\#4.1}.} Whether the gap
is caused by long-\glsterm{chain-of-thought}{CoT} SFT, by \glsterm{rlvr-3}{RLVR} post-training, or by general scale
is not separated by the published 2026 ablations. The hedge that
\citet[\S 6]{wei2022-cot} themselves placed on the cognitive
interpretation of \glsterm{chain-of-thought}{CoT} --- ``\glsterm{chain-of-thought}{chain-of-thought} does not necessarily
reflect model reasoning in any cognitive sense'' --- has aged into the
load-bearing question of \cref{sec:cot-faithfulness}, which catalogues
the 2025--2026 evidence and the open ablations needed to settle it.
\end{contradiction}

\subsection{What this section sets up}
\label{sec:cot-forward}

The decoding factorisation \cref{eq:cot-decode} is the substrate every
remaining section of this paper builds on. \Cref{sec:self-consistency}
takes the cheapest leg of the inference-time compute triangle by
sampling \cref{eq:cot-decode} $N$ times at $T > 0$ and majority-voting
the extracted answers, with no verifier or training change; the gain
is the simplest demonstration that the per-sample noise in $z$ is
recoverable when the answer space is low-cardinality.
\Cref{sec:test-time-compute} promotes both the trace length $|z|$ and
the sample count $N$ from prompting choices to a joint compute axis,
transcribes \citet{snell2024}'s compute-optimal selection problem, and
presents the FLOPs-matched headline that on easy and intermediate
prompts \glsterm{test-time-compute}{test-time compute} substitutes for ${\sim}14\times$ more
pretraining FLOPs. \Cref{sec:process-reward-models} adds the verifier
leg by scoring per-step prefixes of $z$, and \cref{sec:rlvr,sec:r1-family}
add the RL leg by training $p_\theta$ so that the trace distribution
itself concentrates on correct answers. The \glsterm{chain-of-thought}{CoT} phenomenon, in
\citet{wei2022-cot}'s 2022 form, was the first leg of this triangle to
appear in the literature; the rest of the paper traces what happens
when the field stops reading it as a prompting trick and starts
treating it as a compute axis with its own \glsterm{scaling-law}{scaling law}.
